\documentclass[../../../include/open-logic-section]{subfiles}

\begin{document}

\olfileid{sth}{cardinals}{cp}
\olsection{కాంటర్ సూత్రం}

\olref[ordinals][vn]{sec}ను మళ్లీ గుర్తుచేసుకోండి. అక్కడ సుక్రమిత సమితుల గురించి చర్చిస్తూ, సుక్రమాలకు ప్రతినిధులుగా నిలిచే వస్తువులు ఉంటే బాగుంటుందని సూచించాం. ఆ ఉద్దేశంతో క్రమసంఖ్యలను ప్రవేశపెట్టి, అవి ఆశించినట్లుగానే పనిచేస్తాయని \olref[ordinals][ordtype]{ordtypesworklikeyouwant}లో చూపాం, అంటే:
\[
	\ordtype{A, <} = \ordtype{B, \lessdot} 
	\text{ అప్పుడూ అప్పుడే } \tuple{A, <} \isomorphic \tuple{B, \lessdot}.
\]
ఇంకా వెనక్కి, \olref[sfr][siz][equ]{sec}ను గుర్తుచేసుకోండి. అక్కడ అనౌపచారికంగా పనిచేస్తూ సమితి ``పరిమాణం'' అనే భావనను ప్రవేశపెట్టాం. ముఖ్యంగా, $f \colon A \to
B$ అనే \tetoken{ద్వైజెక్టివ్ ప్రమేయం}{!!a{bijection}} ఉన్నప్పుడూ అప్పుడే రెండు సమితులు సమసంఖ్యకాలు, $\cardeq{A}{B}$, అని చెప్పాం. సుక్రమం కంటే ఇది స్వతః \emph{సరళమైన} భావన: మన ఆసక్తి \tetoken{ద్వైజెక్టివ్ ప్రమేయాల}{!!{bijection}s}పైనే; క్రమ-సమరూపతలలో వలె ఆ \tetoken{ద్వైజెక్టివ్ ప్రమేయాలు}{!!{bijection}s} ``ఏ నిర్మాణాన్నైనా పరిరక్షిస్తాయా'' అనే విషయం ఇక్కడ ముఖ్యం కాదు.

ఇది సహజమైన ఆలోచనకు దారి తీస్తుంది. సుక్రమాలను కొలిచేందుకు కొన్ని వస్తువులను, \emph{క్రమసంఖ్యలను}, ప్రవేశపెట్టినట్లే, పరిమాణాన్ని కొలిచేందుకు కొన్ని వస్తువులను, \emph{కార్డినల్ సంఖ్యలను}, ప్రవేశపెట్టవచ్చు. అదే ఈ అధ్యాయం లక్ష్యం.

ఈ కార్డినల్ సంఖ్యలు ఏమిటో చెప్పే ముందు, అవి పాటించాల్సిన సూత్రాన్ని నిర్దేశిద్దాం. $X$ సమితి కార్డినాలిటీని $\card{X}$గా రాస్తే, అవి ఈ నియమాన్ని పాటించాలని కోరుకుంటాం:
\[
	\card{A} = \card{B} \text{ అప్పుడూ అప్పుడే } \cardeq{A}{B}.
\]
దీన్ని \emph{కాంటర్} సూత్రం అంటాం; బహుశా కాంటరే దీన్ని మొదట స్పష్టంగా మనసులో ఉంచుకున్నారు. (\olref[hp]{sec}లో దీనికి \emph{హ్యూమ్} సూత్రంతో ఉన్న సంబంధం గురించి మరింత చెబుతాం.) కాబట్టి ప్రతి $X$కు $\card{X}$ను కాంటర్ సూత్రాన్ని నెరవేర్చేలా నిర్వచించడమే మన లక్ష్యం.

\end{document}
